\section{Discussion}\label{sec:discussion}

\subsection{What the size of the effect depends on}

The ten point estimates in Table~\ref{tab:main} differ across libraries and backbones. The
library medians of exposure are similar, but class-level exposures and other properties vary.
We did not fit an explanatory model or establish a variance decomposition, so we cannot determine
which factors explain the observed spread.

Headroom constrains the scale. A library whose control arm already reaches 0.95 cannot show a large
absolute gain, whatever the split does; CTIS loses at most 4.6 points to Swin-T and 9.2 to
AlexNet. In bel canto the recording-level control reaches the low 0.9s and the exposure arm
saturates at 1.00 for most seeds, so its recording-level contrast is limited by a ceiling and its
interval is wide for that reason as well as for the small panel of 40 test recordings.
Rescaling by headroom, as Table~\ref{tab:headroom} does, reduces but does not remove
the dispersion (0.219--0.873, still about four-fold): the normalised contrast is 22--37\% of the score headroom in CTIS, 38--48\% in
CNPM, 47--68\% in bel canto, 49--68\% in music genre, and 83--87\% in pianos. Absolute points
answer the question ``how much would the leaderboard move''; the rescaled number answers the
question ``how large is the contrast relative to the baseline's remaining score headroom''.
Because $1-\text{macro-F1}$ is not an error count, this is a descriptive normalised
effect size, not a division of predictions into memorisation and transfer.

The exposure ratio mainly reflects the chosen test share, half-pool allocation and training
budget, giving approximately 0.167. The five library medians range from 0.160 to 0.170, while
class-level values vary. These summaries alone cannot establish the role of exposure in the
between-library differences; task difficulty, class structure and recording similarity may also
matter, and our measurements do not isolate their contributions.

\subsection{Interpreting the primary result}

The internal plan identifies CTIS as the primary library because it is the largest and its class
structure supports the main design. There is no independent public timestamp establishing when
that choice was fixed. CTIS has the smallest observed absolute effects, with limited score
headroom for both backbones. Conditional intervals exclude zero, but only the Swin-T interval is
wholly below the study-specific three-point SESOI. These results distinguish statistical evidence
from practical importance; they do not establish the timing of the internal choice.

The four additional libraries show larger observed changes, from $+8.4$ to $+28.7$ points,
compared with $+2.02$ (AlexNet) and $+1.70$ (Swin-T) in CTIS. The tasks and class structures
differ, so these comparisons are descriptive.

The C1 control sharpens what the primary contrast can and cannot mean. It supplies the same volume
of training slices from non-test content, and the exposure arm still leads it by $+1.96$ points
under AlexNet ($+0.74$ to $+2.99$) and $+1.07$ under Swin-T ($+0.23$ to $+1.96$), so the
difference between the two arms is not explained by training volume. What survives is a contrast
that also changes the provenance of the injected content and the diversity of its recordings, which
the design does not separate: the A1--A0 difference is consistent with an advantage of
within-recording exposure, and it does not isolate one.

\subsection{What a reporting standard should require}

Four reporting practices would make evaluation protocols easier to interpret.

\emph{Name the split unit.} ``Randomly split 8:1:1'' is ambiguous only because the release format
makes slices look atomic. A statement of the form ``split by recording, 8:1:1 over recordings,
stratified by class'' identifies the intended unit; releasing the split assignments and grouping
rules allows the implementation to be checked.

\emph{Report the slice-per-recording distribution.} The distribution of $K$, not its mean alone,
determines the exposure probabilities under a specified slice-level split; the mean ranges from 8.7 to 48.9 across these
five libraries, so a single pooled statement about ``the dataset'' would be misleading even if
the split unit were named.

\emph{Report capacity when a stricter split is claimed.} If a paper says it splits by recording,
the reader needs to know whether that was feasible at the stated class counts: $q_c$,
$m_c$, and the resulting $\alpha_c$ document the constraints. CNPM and bel canto retain only
54 and 40 test recordings, respectively, so inference at recording level is especially uncertain. Our feasibility table is released with the splits precisely so that this
arithmetic can be checked rather than trusted.

\emph{Preserve the evaluation record.} Evaluating the held-out panel after training without
selecting on it limits selection bias, and stored predictions permit re-analysis. Our local
staging and sealing records do not independently establish when choices were made.

\subsection{Relation to prior audits}

Our result is consistent with Sturm~\cite{sturm2012gtzan} in kind and different in object: Sturm found
that a benchmark's content --- its labels and its duplicates --- did not support the conclusions drawn from it; we
show how one reconstruction of an ambiguously described evaluation recipe can affect scores.
The original realised split is unknown, so we cannot attribute the observed score change to that implementation.
The distinction matters for repair. Mislabelled audio has to be re-annotated; a
slice-level split can be rebuilt from recording identifiers if class capacity permits it, as
we demonstrate for the retained classes in these five libraries. Split construction itself is
CPU work; training costs are itemised separately in Section~\ref{sec:results}.

\subsection{Limitations and threats to validity}

\emph{One feature family.} We audited MEL only. The framework also releases chroma and CQT, and
the leakage mechanism is identical for them, but the magnitude need not be: chroma discards
timbre and may carry less recording identity.

\emph{What the contrast identifies.} A1 holds the training budget and the C-side recording
count fixed, but it does add pool recordings, so its training set contains more distinct
recordings than A0's. The contrast is therefore the effect of equal replacement under the
corpus's own composition --- the rule a reader applying this protocol would follow --- and not a
decomposition of memory into a content component and a diversity component. C1 addresses this
partly: equal-volume non-test injection changes macro-F1 by $+0.05$ and $+0.62$ points compared
with $+2.02$ and $+1.70$ for A1. C1 also changes recording diversity and uses R content; the
C1--A1 difference therefore cannot isolate a pure memorisation component.

\emph{Two backbones, three seeds.} The primary analysis uses AlexNet and Swin-T at seeds
0--2. The primary intervals condition on these runs; the supplementary joint bootstrap exposes
sensitivity to their dispersion, but three seeds cannot characterise the distribution of training
outcomes. We claim neither upper nor lower bounds for other architectures or training recipes.

\emph{One split realisation.} The recording-level split has a single seed. Capacity numbers are
computed for that realisation, and per-class exposure is reported rather than tested, because
the capacity table stores the mean of 400 random draws while a split realises one.

\emph{Reconstruction of the audited protocol.} We could not retrieve the framework's own split
files; the A2 arm reconstructs one slice-level interpretation of the stated 8:1:1 protocol and
may differ from the authors' actual choice of split unit, seed and stratification. The A1$-$A0
contrast does not depend on that reconstruction and is the primary contrast in the internal plan.

\emph{Version of record.} We checked the published split description in TISMIR 8(1),
pp. 22--38 (Section~\ref{sec:framework}). The original realised splits remain unavailable, so
our contrasts cannot be attributed numerically to the original authors' models.

\emph{External validity.} All five libraries are Chinese-music collections. Nothing in the
mechanism is language-specific. In the single-slice controls recording- and slice-level units
coincide by construction; the low-exposure control only approximates this condition. These
sanity checks are consistent with, but do not by themselves establish, slice multiplicity as the
operative explanation. The control libraries differ in task and classes, and one contrast excludes
zero. We do not claim that the magnitudes transfer to other benchmark suites without measurement.
